\sezione{9}{Esercizio completo: la biblioteca}{slide 45–47 e 153–157}

\begin{inbreve}
Un esercizio d'esame (13/09/2002) che usa quasi tutti i costrutti: entità distinte per concetti diversi, relationship uno a molti, una \hl{catena di identificatori esterni}, un attributo di relationship e una prima ridondanza. Prova a risolverlo da solo prima di leggere la soluzione.
\end{inbreve}

\subsection{Le specifiche}

\begin{tecnica}[Testo dell'esercizio]
Definire uno schema E-R per una biblioteca, con le seguenti specifiche:
\begin{itemize}
  \item oggetto dei prestiti sono esemplari (detti anche copie) di singoli volumi, identificati attraverso un numero di inventario; ogni volume è relativo a una specifica edizione (che può essere articolata in più volumi, anche in modo diverso dalle altre edizioni) di un'opera;
  \item un volume può essere presente in più copie;
  \item un'edizione è caratterizzata dall'opera, dalla collana e dall'anno;
  \item riassumendo ed esemplificando, è possibile prendere in prestito la seconda copia del terzo volume de «I Miserabili», edizione Mondadori, collana Oscar, del 1975;
  \item ogni collana ha un nome, un codice e un editore; ogni editore ha un nome e un codice;
  \item ogni opera ha un titolo, un autore e un anno di prima pubblicazione;
  \item per ogni prestito in corso (quelli conclusi non interessano) sono rilevanti la data prevista di restituzione e l'utente (che può avere più volumi in prestito contemporaneamente), con codice identificativo, nome, cognome e recapito telefonico.
\end{itemize}
\end{tecnica}

\begin{consiglio}[Come affrontarlo]
\begin{passi}
  \item Sottolinea i \hl{sostantivi} che hanno proprietà proprie: sono candidati a diventare entità (opera, edizione, volume, copia, collana, editore, utente).
  \item Cerca i \hl{legami} e per ognuno chiediti le cardinalità, mettendoti nei panni di una singola occorrenza.
  \item Per ogni entità chiediti \hl{che cosa la identifica}: se la numerazione «riparte» (il terzo volume \emph{di un'edizione}), serve un identificatore esterno.
  \item Verifica lo schema con l'\hl{esempio} delle specifiche.
\end{passi}
\end{consiglio}

\subsection{La soluzione}

\par\noindent\begin{minipage}{\linewidth}
\begin{schema}[x=1cm,y=1cm,scale=.93,every node/.append style={transform shape}]
  \node[ent] (Ed) at (0,0) {Editore};
  \node[ent] (V) at (6,0) {Volume};
  \node[ent] (Cp) at (11,0) {Copia};
  \node[ent] (Co) at (0,-3.2) {Collana};
  \node[ent] (E) at (6,-3.2) {Edizione};
  \node[ent] (U) at (11,-3.2) {Utente};
  \node[ent] (O) at (6,-6.4) {Opera};
  \node[rel,minimum width=1.4cm] (r1) at (0,-1.6) {\phantom{x}};
  \node[rel,minimum width=1.4cm] (r2) at (3,-3.2) {\phantom{x}};
  \node[rel,minimum width=1.4cm] (r3) at (6,-1.6) {\phantom{x}};
  \node[rel,minimum width=1.4cm] (r4) at (8.5,0) {\phantom{x}};
  \node[rel] (Pr) at (11,-1.6) {Prestito};
  \node[rel,minimum width=1.4cm] (r5) at (6,-4.8) {\phantom{x}};
  \draw[lk] (Ed) -- (r1) node[card,pos=.4,right]{(0,N)};
  \draw[lk] (r1) -- (Co) node[card,pos=.6,right]{(1,1)};
  \draw[lk] (Co) -- (r2) node[card,pos=.35,above]{(0,N)};
  \draw[lk] (r2) -- (E) node[card,pos=.25,above]{(1,1)};
  \draw[lk] (V) -- (r3) node[card,pos=.15,right]{(1,1)};
  \draw[lk] (r3) -- (E) node[card,pos=.6,right]{(0,N)};
  \draw[lk] (V) -- (r4) node[card,pos=.35,above]{(0,N)};
  \draw[lk] (r4) -- (Cp) node[card,pos=.6,above]{(1,1)};
  \draw[lk] (Cp) -- (Pr) node[card,pos=.4,right]{(0,1)};
  \draw[lk] (Pr) -- (U) node[card,pos=.6,right]{(0,N)};
  \draw[lk] (E) -- (r5) node[card,pos=.8,right]{(1,1)};
  \draw[lk] (r5) -- (O) node[card,pos=.6,right]{(0,N)};
  % attributi
  \idattributo{(-.5,.41)}{(-.5,1.3)}{Codice}{above}
  \attributo{(.5,.41)}{(.5,1.3)}{Nome}{above}
  \idattributo{(-.5,-3.61)}{(-.5,-4.5)}{Codice}{below}
  \attributo{(.5,-3.61)}{(.5,-4.5)}{Nome}{below}
  \attributo{(V.west)}{(3.5,0)}{Numero}{above}
  \idattributo{(Cp.north)}{(11,1.2)}{NumInv}{above}
  \attributo{(Pr.east)}{(12.9,-1.6)}{Data}{right}
  \attributo{(E.east)++(0,.15)}{(8.3,-2.75)}{NVolumi}{right}
  \attributo{(E.east)++(0,-.15)}{(8.3,-3.75)}{Anno}{right}
  \idattributo{(U.south west)++(.3,0)}{(9.9,-4.4)}{Codice}{below}
  \attributo{(U.south)}{(11,-4.4)}{Cognome}{below}
  \attributo{(U.south east)++(-.3,0)}{(12.1,-4.4)}{Nome}{below}
  \attributo{(U.east)}{(13,-3.2)}{Telefono}{right}
  \idattributo{(O.west)}{(3.9,-6.4)}{Codice}{left}
  \attributo{(O.east)++(0,.2)}{(8.2,-5.95)}{Titolo}{right}
  \attributo{(O.east)}{(8.2,-6.4)}{Autore}{right}
  \attributo{(O.east)++(0,-.2)}{(8.2,-6.85)}{Anno}{right}
  % identificatore esterno di Volume: Numero + Edizione
  \draw[idl] (4.2,.5) -- (4.2,-1.0) -- (6.35,-1.0); \node[iddot] at (4.2,.5) {};
  % identificatore esterno di Edizione: Collana + Opera + Anno
  \draw[idl] (4.35,-2.8) -- (4.35,-3.95) -- (7.9,-3.95) -- (7.9,-3.55); \node[iddot] at (4.35,-2.8) {};
\end{schema}
\didascalia{Soluzione (slide 155). Le relationship senza nome sono quelle disegnate senza nome nelle slide. In alternativa, e forse meglio, l'edizione può avere un proprio codice (l'ISBN).}
\end{minipage}\par

\begin{tecnica}[Le scelte, una per una]
\begin{itemize}
  \item \hl{Opera, Edizione, Volume, Copia} sono quattro entità: hanno occorrenze diverse (capitolo 2).
  \item \hl{Copia} è identificata dal numero di inventario (NumInv), come dicono le specifiche.
  \item \hl{Volume} è identificato da Numero \key{e} Edizione: «il terzo volume» è tale solo all'interno della sua edizione.
  \item \hl{Edizione} è identificata da Collana, Opera e Anno («caratterizzata dall'opera, dalla collana e dall'anno»): entrambe le relationship hanno $(1,1)$ dal lato di Edizione.
  \item \hl{Prestito} è una relationship fra Copia e Utente con l'attributo \emph{Data} (prevista di restituzione). Interessano solo i prestiti \key{in corso}: ogni copia ha al più un prestito, da cui $(0,1)$; un utente può avere più prestiti, da cui $(0,N)$.
  \item \hl{Editore} si raggiunge attraverso la Collana: non serve collegarlo anche all'Edizione.
\end{itemize}
\end{tecnica}

\begin{esempio}[Verifica con l'esempio delle specifiche]
«La seconda copia del terzo volume de \emph{I Miserabili}, edizione Mondadori, collana Oscar, del 1975»: la \key{Copia} (numero di inventario) → il \key{Volume} 3 → l'\key{Edizione} (collana Oscar, opera \emph{I Miserabili}, anno 1975) → la \key{Collana} Oscar → l'\key{Editore} Mondadori. Ogni passaggio dell'esempio corrisponde a una relationship dello schema: è la \hl{catena di identificatori esterni} (Volume → Edizione → Collana e Opera), come per le coordinate bancarie.
\end{esempio}

\subsection{Errori frequenti}

\begin{attenzione}[Biblioteca: errori frequenti (slide 156)]
\begin{itemize}
  \item \hl{Un'unica entità «Libro»}: opera, edizione, volume e copia sono concetti diversi.
  \item \hl{Volume identificato dal solo numero}: serve l'identificatore esterno verso l'edizione.
  \item \hl{Editore collegato anche all'edizione}: è ridondante, si raggiunge attraverso la collana.
  \item \hl{Data prevista di restituzione su Copia o Utente}: è una proprietà del prestito.
  \item \hl{Prestiti rappresentati come storico}: interessano solo quelli in corso, quindi $(0,1)$ dal lato di Copia.
\end{itemize}
\end{attenzione}

\begin{mlbox}[Aspetti da sottolineare (slide 157)]
\begin{itemize}
  \item \key{Codici introdotti dal progettista}: per l'opera, identificarla con autore e titolo sarebbe «fragile»; per l'edizione, in alternativa, l'ISBN.
  \item \key{Gli autori sono semplici valori}: va bene perché non ci viene chiesto altro, e così si gestisce anche il caso di più autori (un valore che li contiene).
  \item \key{NVolumi è derivabile} contando i volumi dell'edizione: è un primo esempio di \hl{ridondanza}, cioè di informazione che si può ricavare da altre.
\end{itemize}
\end{mlbox}

\begin{esercizi}
\begin{passi}
  \item Modifica lo schema per conservare anche lo \emph{storico} dei prestiti conclusi, con data di inizio e data di restituzione effettiva.
  \item Come cambia lo schema se l'edizione ha un codice ISBN?
\end{passi}
\end{esercizi}

\begin{appunti}{Riepilogo del capitolo}
\pag{47} Quattro entità per il «libro»: Opera, Edizione, Volume, Copia.
\pag{155} Catena di identificatori esterni: Volume → Edizione → Collana e Opera.
\pag{156} Errori: Libro unico, Volume senza identificatore esterno, Editore ridondante, data sul posto sbagliato, storico non richiesto.
\pag{157} Codici introdotti dal progettista; autori come valori; NVolumi ridondante.
\end{appunti}
